% Arquivo: latex/1-introducao.tex
% Seção 1: Introdução

A desigualdade regional constitui um dos traços estruturais mais persistentes da economia brasileira. As regiões Norte e Nordeste concentram historicamente os piores indicadores de renda, educação e infraestrutura, enquanto Sul e Sudeste respondem pela maior parcela do produto nacional \cite{Ferreira1995,Silveira2011}. Embora estudos identifiquem processo de convergência dos níveis regionais de renda ao longo das últimas décadas, a velocidade é insuficiente para eliminar as disparidades em horizonte razoável: \citeonline{Cruz2014} estimam que, mantido o ritmo de crescimento diferencial observado, o Nordeste levaria cinquenta anos para atingir 75\% do PIB \textit{per capita} nacional. Dados recentes indicam que cerca de 90\% dos municípios das áreas de atuação da SUDAM, SUDENE e SUDECO classificados na faixa inferior de renda em 2002 permaneciam na mesma condição em 2021, evidenciando elevada persistência do baixo desenvolvimento \cite{Portugal2024}.

Nesse contexto, a Política Nacional de Desenvolvimento Regional (PNDR), instituída pelo Decreto nº 6.047/2007 e reformulada em 2019 e 2024, opera por meio de três instrumentos de financiamento voltados às regiões Norte, Nordeste e Centro-Oeste: os Fundos Constitucionais de Financiamento (FNE, FNO e FCO), que oferecem crédito subsidiado aos setores produtivos locais; os Fundos de Desenvolvimento Regional (FDNE, FDA e FDCO), direcionados a projetos estruturantes de grande porte; e os Incentivos Fiscais administrados pela SUDAM e SUDENE, que operam via renúncia tributária para atrair investimentos privados \cite{Portugal2020}. O volume de recursos públicos mobilizado por esses instrumentos é expressivo: no período 2002--2021, as aplicações acumuladas dos Fundos Constitucionais totalizaram cerca de R\$ 549 bilhões (valores de 2020), os Fundos de Desenvolvimento liberaram aproximadamente R\$ 18,9 bilhões em financiamentos, e os Incentivos Fiscais da SUDAM e SUDENE responderam por R\$ 126 bilhões em gastos tributários entre 2015 e 2024 \cite{Portugal2024}.

Apesar da longevidade dos programas --- os Fundos Constitucionais operam ininterruptamente desde 1989 e os Incentivos Fiscais desde a década de 1960 --- e do volume substancial de recursos envolvidos, a persistência das disparidades regionais torna imprescindível a avaliação rigorosa da eficácia desses instrumentos. A adoção de práticas sistemáticas de monitoramento e avaliação é condição necessária para verificar se os objetivos da política estão sendo alcançados, identificar falhas de desenho e fundamentar ajustes na alocação dos recursos \cite{Lima2025}. No caso de políticas de desenvolvimento regional, a avaliação enfrenta desafios metodológicos específicos: a dificuldade de construir grupos contrafactuais adequados, a multiplicidade de indicadores de resultado, a presença de efeitos de transbordamento espacial e a defasagem temporal entre a intervenção e a maturação dos seus efeitos \cite{Oecd2025}.

Essas dificuldades fazem com que métodos experimentais randomizados sejam raramente factíveis na avaliação de políticas regionais, levando os pesquisadores a recorrer a métodos de painel com efeitos fixos, modelos espaciais, técnicas de \textit{matching} e desenhos quase experimentais como diferenças em diferenças, regressão com descontinuidade e controle sintético. Nas últimas duas décadas, uma literatura empírica crescente buscou estimar os efeitos dos instrumentos da PNDR sobre indicadores socioeconômicos locais empregando essas abordagens. Os estudos, contudo, encontram-se dispersos em periódicos nacionais e internacionais, \textit{working papers}, textos para discussão e anais de congressos, dificultando uma visão integrada do estado das evidências.

Revisões prévias da literatura puderam identificar padrões importantes nessa produção, embora com escopos parciais. \citeonline{Resende2013} realizou levantamento da literatura de avaliação de políticas regionais no Brasil e apontou relativa escassez de estudos empíricos no período 2000--2009, com apenas duas avaliações de impacto identificadas. \citeonline{Portugal2017} analisou estudos realizados entre 2005 e 2016 e constatou que os Fundos Constitucionais são os instrumentos mais investigados, com evidências de impacto positivo sobre emprego e massa salarial, mas com resultados divergentes sobre o PIB \textit{per capita}. \citeonline{CarneiroCambota2018} sintetizaram evidências empíricas específicas do FNE, concluindo por efeitos geralmente positivos, porém heterogêneos entre setores e mais pronunciados para micro e pequenas empresas. \citeonline{Shirasu2023} revisaram estudos sobre incentivos fiscais, mas relataram pouca aplicação ao contexto específico da PNDR. Nenhuma dessas revisões, contudo, adotou protocolo sistemático de busca e seleção, tampouco abrangeu simultaneamente os três instrumentos em horizonte temporal amplo.

A pergunta de pesquisa que orienta este artigo é: qual é o estado das evidências empíricas sobre os efeitos dos Fundos Constitucionais, Fundos de Desenvolvimento e Incentivos Fiscais da PNDR sobre indicadores socioeconômicos locais --- em particular PIB \textit{per capita}, emprego formal e salário médio --- no período de 2005 a 2026? O artigo contribui para preencher a lacuna identificada ao conduzir uma revisão sistemática conforme o protocolo PRISMA 2020 \cite{Page2021}, com busca em cinco bases de dados acadêmicas (Scopus, SciELO, Catálogo de Teses e Dissertações da CAPES, EconPapers/RePEc e Encontros da ANPEC) e seleção de 34 estudos empíricos.

Para além da síntese comparativa dos resultados por instrumento e abordagem metodológica, o artigo incorpora inovações procedimentais que o diferenciam das revisões anteriores: (i)~processo de deduplicação em quatro fases (DOI exato, título por similaridade \textit{fuzzy}, PDF idêntico e verificação manual de textos para discussão e versões de congresso); (ii)~classificação assistida por modelo de linguagem (Gemini 2.0 Flash), em três estágios de análise, com revisão manual subsequente; e (iii)~cálculo de índice de citação cruzada entre os estudos selecionados, que permite identificar os trabalhos mais influentes dentro do campo e mapear a rede de referências da literatura.

Os principais resultados da revisão indicam: (i)~expressiva assimetria de evidências, com 28 dos 34 estudos concentrados nos Fundos Constitucionais e predominância do FNE; (ii)~maior convergência nas estimativas de impacto sobre o emprego formal do que sobre o PIB \textit{per capita}; (iii)~eficácia dos Fundos Constitucionais condicionada ao nível inicial de renda dos municípios, com efeitos positivos em faixas intermediárias de desenvolvimento e nulos em municípios de baixa renda; (iv)~indícios de impacto expressivo dos Fundos de Desenvolvimento sobre o PIB \textit{per capita} (entre 19\% e 24\%), porém com base empírica ainda restrita; e (v)~ausência de ganhos de produtividade associados à geração de empregos nos instrumentos avaliados.

O restante do artigo organiza-se da seguinte forma: a Seção~\ref{sec:politica_regional} contextualiza a política regional no Brasil e descreve o funcionamento dos três instrumentos; a Seção~\ref{sec:metodo} apresenta a estratégia metodológica da revisão sistemática; a Seção~\ref{sec:resultados} discute os resultados organizados por instrumento; e a Seção~\ref{sec:consideracoes} apresenta as considerações finais, incluindo implicações para o desenho da política e lacunas prioritárias para a agenda de pesquisa.
